\documentclass[11pt]{article}
    \usepackage[utf8]{inputenc}
    \usepackage[T1]{fontenc}
    \usepackage[margin=1in]{geometry}
    \usepackage{hyperref}
    \usepackage{enumitem}
    \title{What is Coding?}
    \author{Piter Garcia}
    \date{March 20, 2026}
    \begin{document}
    \maketitle
    \section*{Quick Scan}
    \begin{itemize}[leftmargin=*]
    \item Grade: Grade 4
    \item Workbook sessions: Session 22
    \item Class blocks: Day 1 1:15-2:10 Baris, Day 2 1:15-2:10 Fowler, Day 3 1:15-2:10 Schrank, Day 4 1:15-2:10 Autore
    \item Reference file: what-is-coding.bib
    \end{itemize}
    \section*{School Planning Goals and Inclusion}
    \begin{itemize}[leftmargin=*]
    \item What is expected: - I can decompose a task into smaller named tasks (4-6.CT.4) | - I can reframe a sequence of steps as an encoded program (4-6.CT.9)
    \item What we achieved: No direct transcript match is attached yet. Treat this lesson as workbook-backed and enrich it when a matching classroom transcript is available.
    \item What needs improvement: stronger proof, cleaner assessment notes, and tighter lesson artifacts.
    \item How to improve: add transcript proof, one saved artifact, and clearer success criteria.
    \item Inclusion focus: use visual modeling, chunked directions, predictable routines, and flexible participation roles.
    \end{itemize}
    \section*{Official References}
    \begin{itemize}[leftmargin=*]
    \item \url{https://csteachers.org/k12standards/interactive/}
\item \url{https://code.org/curriculum/elementary-school}

    \end{itemize}
    \end{document}
